\begin{helper}
Fix a terminal coordinate set \(U\).  Let
\(\mathsf{Realized}(\beta,\tau)\) mean that the branch/transcript cell
\((\beta,\tau)\) is nonempty, and let
\(\mathsf{Comp}(A,\beta,\tau)\) mean that a complete terminal assignment
\(A\subseteq U\) is compatible with that realized cell.  Suppose also that
\(\mathsf{Sibling}(A,\beta,\tau,e)\) denotes the selected-present sibling
created by changing a selected absence \(e\).

Assume the selected-sibling pruning rule: if
\(\mathsf{Sibling}(A,\beta,\tau,e)\) occurs for a realized cell, then the
same assignment \(A\) is compatible with no realized cell at all.  Assume
also the functional scan rule: whenever \(A\subseteq U\) is compatible with
a realized cell \((\beta,\tau)\), the branch label is the value selected
from \(A\)'s complete terminal blue trace and the deterministic prefix scan
on \(A\) returns \(\tau\).

Then one terminal assignment is compatible with at most one realized
branch/transcript cell.  More precisely, if \(A\subseteq U\) is compatible
with realized cells \((\beta,\tau)\) and \((\beta',\tau')\), then no
selected-present sibling of either cell is active at \(A\), and
\[
\beta=\beta',
\qquad
\tau=\tau'.
\]
\end{helper}
\begin{proof}
Fix \(A\subseteq U\) and suppose that \(A\) is compatible with the two
realized cells \((\beta,\tau)\) and \((\beta',\tau')\).  First check the
selected-sibling cases.  If
\(\mathsf{Sibling}(A,\beta,\tau,e)\) held for some selected coordinate
\(e\), the pruning hypothesis for the realized cell \((\beta,\tau)\) would
say that \(A\) is compatible with no realized cell.  This contradicts the
assumed compatibility of \(A\) with \((\beta,\tau)\).  Therefore no
selected-present sibling is active for the first cell.  The same argument,
with \((\beta',\tau')\) in place of \((\beta,\tau)\), shows that no
selected-present sibling is active for the second cell.

The hypotheses of
\(\noderef{FixedSetHistoryCellRedFunctionalScanDisjointness}\) are now
available for the two compatible realized cells: \(A\subseteq U\), both
cells are realized, both compatibility relations hold, and the branch and
scan outputs are functional from \(A\).  Applying that node to the first
cell gives
\[
\beta=\operatorname{branchOfAssignment}(A),
\qquad
\operatorname{scanTranscript}(A)=\tau.
\]
Applying it to the second cell gives
\[
\beta'=\operatorname{branchOfAssignment}(A),
\qquad
\operatorname{scanTranscript}(A)=\tau'.
\]
The two branch labels are therefore equal to the same selected branch, and
the two transcripts are the same deterministic scan output.  Hence
\(\beta=\beta'\) and \(\tau=\tau'\).
\end{proof}
